% package
\usepackage{alltt}
\usepackage{amsfonts}
\usepackage{amsmath}
\usepackage{hyphenat}
\usepackage{amssymb}
\usepackage{amsthm}
\usepackage{bm}
\usepackage{booktabs}
\usepackage{caption}
\usepackage{color}
\usepackage[colorlinks]{hyperref}
\usepackage{enumitem}
\usepackage{etoolbox}
\usepackage{fancyhdr}
\usepackage{flafter}
\usepackage{geometry}
\usepackage{graphicx}
\usepackage{mathalfa}
\usepackage{mathdots}
\usepackage{mathrsfs}
\usepackage{mathtools}
% \usepackage{microtype}
\usepackage{multirow}
\usepackage{multicol}
\usepackage{pdflscape}
\usepackage{pgffor}
\usepackage{pgfplots}
\usepackage{siunitx}
\usepackage{slashed}
\usepackage{subfigure}
\usepackage{tabularx}
\usepackage{tikz}
\usepackage{tkz-euclide}
\usepackage[normalem]{ulem}
\usepackage{xcolor}
\usepackage{xparse}
\usepackage{xpatch}
\usepackage[all]{xy}
\usepackage{imakeidx}

% layout
\pagestyle{fancy}
\fancyhf{}
\fancyfoot[C]{\thepage}
\lhead{\emph{\nouppercase{\leftmark}}}
\setlength{\headheight}{14.5pt}

% enumeration
\renewcommand{\theenumi}{\alph{enumi}}
\renewcommand\labelenumi{(\theenumi)}

% title
\author{ioshift}
\date{\today}
\title{\Huge{\bf{\ntitle}}}

% Others
\pgfplotsset{compat=1.18}
\setlength{\parindent}{0em}
\setlist[itemize,2]{label=$\circ$}


%%%%%%%%%%%%%%%%%%%%%%%%%%
%%%%%%%%%% math %%%%%%%%%%
%%%%%%%%%%%%%%%%%%%%%%%%%%

% Theorems
\numberwithin{equation}{subsection}
\theoremstyle{definition}
\newtheorem{thm}{Theorem}[subsection]
\renewcommand{\thethm}{%
  \ifnum\value{subsection}=0
    \thesection.\arabic{thm}%
  \else
    \thesubsection.\arabic{thm}%
  \fi
}
\newtheorem{axiom}[thm]{Axiom}
\newtheorem{conjecture}[thm]{Conjecture}
\newtheorem{cor}[thm]{Corollary}
\newtheorem{eg}[thm]{Example}
\newtheorem{prop}[thm]{Proposition}
\newtheorem{lemma}[thm]{Lemma}
% \newtheorem{prob}{Problem}[section]
\newtheorem{ex}[thm]{Exercise}

\newtheorem*{aim}{Aim}
\newtheorem*{assumption}{Assumption}
\newtheorem*{claim}{Claim}
\newtheorem*{Defi}{Definition}
\newtheorem*{fact}{Fact}
\newtheorem*{law}{Law}
\newtheorem*{remark}{Remark}
\newtheorem*{note}{Note}
\newtheorem*{rrule}{Rule}
\newtheorem*{Prop}{Proposition}
\newtheorem*{question}{Question}
\newtheorem*{warning}{Warning}

\renewcommand{\proofname}{\textbf{Proof}}
\newenvironment{sol}{\begin{proof}[Solution]}{\end{proof}}
\newenvironment{counter}{\begin{proof}[\textbf{Counterexample}]}{\end{proof}}

% Special sets
\newcommand{\C}{\mathbb{C}}
\newcommand{\R}{\mathbb{R}}
\newcommand{\Q}{\mathbb{Q}}
\newcommand{\Z}{\mathbb{Z}}
\newcommand{\N}{\mathbb{N}}
\newcommand{\F}{\mathbb{F}}

\newcommand{\topo}{\mathcal{T}}
\newcommand{\B}{\mathcal{B}}
\newcommand{\cont}{\mathcal{C}}
\renewcommand{\S}{\mathcal{S}}
% Symbols
\newcommand{\noteq}{\not=}
\newcommand{\re}{\mathrm{Re\,}}
\newcommand{\im}{\mathrm{Im\,}}
\newcommand{\nexist}{\not\!\exists\,}
\newcommand{\abs}[1]{\left\lvert #1\right\rvert}
\newcommand{\norm}[1]{\left\lVert #1\right\rVert}
\newcommand{\floor}[1]{\left\lfloor #1\right\rfloor}
\newcommand{\ceil}[1]{\left\lceil #1\right\rceil}
\newcommand{\set}[1]{\left\{#1\right\}}
\newcommand{\closure}[2][3]{%
{}\mkern#1mu\overline{\mkern-#1mu#2}}
\newcommand{\interior}{\mathrm{Int\,}}
\newcommand{\upint}{\mathchoice%
    {\mkern13mu\overline{\vphantom{\intop}\mkern7mu}\mkern-20mu}%
    {\mkern7mu\overline{\vphantom{\intop}\mkern7mu}\mkern-14mu}%
    {\mkern7mu\overline{\vphantom{\intop}\mkern7mu}\mkern-14mu}%
    {\mkern7mu\overline{\vphantom{\intop}\mkern7mu}\mkern-14mu}%
  \int}
\newcommand{\lowint}{\mkern3mu\underline{\vphantom{\intop}\mkern7mu}\mkern-10mu%
\int}
\newcommand{\aand}{\quad\text{and}\quad}
\newcommand{\res}[1]{\!\mid_{#1}}
\newcommand{\normal}{\unlhd}
\renewcommand{\tt}[1]{\quad\text{#1}\quad}
\renewcommand{\d}{\mathop{}\!\mathrm{d}} 
\renewcommand{\emptyset}{\varnothing}
\DeclareMathOperator{\lcm}{lcm}
\DeclareMathOperator{\diam}{diam}
\DeclareMathOperator{\id}{id}
\DeclareMathOperator{\Bd}{Bd}
\DeclareMathOperator{\card}{card}
\DeclareMathOperator{\supp}{supp}


%%% Automaic Definition Glossary
\newtheorem{definition}[thm]{Definition}
\newtheorem{notations}[thm]{Notation}
\newcounter{cnt}
\newcommand{\newterm}[3]{%
    \global\expandafter\unexpanded\expandafter{
    \expandafter\gdef\csname term#1\endcsname}{#2}%
    \global\expandafter\unexpanded\expandafter{
    \expandafter\gdef\csname cat#1\endcsname}{#3}%
}
\newenvironment{defi}[1]{
    \refstepcounter{cnt}\label{\thecnt}
    \newterm{\arabic{cnt}}{#1}{0}
    \begin{definition}\textbf{#1}
    
}{\end{definition}}
\newenvironment{notation}[1]{
    \refstepcounter{cnt}\label{\thecnt}
    \newterm{\arabic{cnt}}{#1}{0}
    \begin{notations}\textbf{#1}
    
}{\end{notations}}
\newcommand{\addelement}[1]{
    \refstepcounter{cnt}\label{\thecnt}%
    \newterm{\arabic{cnt}}{#1}{0}%
}
\newcommand{\addclass}[1]{
    \refstepcounter{cnt}\label{\thecnt}%
    \newterm{\arabic{cnt}}{#1}{1}%
}
\newcommand{\printglossary}{
    \newpage
    \section{Glossary}
    \foreach \x in {1,...,\arabic{cnt}} {
        \ifnum\csname cat\x\endcsname=0
            \hyperref[\x]{\csname term\x\endcsname}\par\addvspace{0.5em}
        \else
            \ifnum\x>1
                \vspace{1em}
            \fi
            {\large\textbf{\csname term\x\endcsname}\par\addvspace{0.7em}}
        \fi
    }
}